\section{Discussion and limitations}
\label{sec:limitations}
The evaluation separates the target, the validity assumptions and the efficiency measure. A joint region must be judged by simultaneous vector coverage. Marginal validity must be accompanied by subgroup diagnostics when subgroup reliability matters. An efficiency comparison needs both coverage and a shared measure of size. These principles describe how to interpret the experiments; the evidence supports only the specific models, distributions and interventions actually tested.

\paragraph{Method choice under the tested conditions.}
Use a joint score or an explicit multiplicity correction when the whole vector must be covered; naive coordinate calibration failed that target. For predeclared observed groups requiring separate reliability, group calibration fits the tested within-group-preserving shift, subject to sparse-group costs. Training-only scaling improves adaptation and expected iid area here, but its measured high-noise-group miss at ten components rules out a general group guarantee. Known-ratio weighting addresses shifted marginal coverage under its assumptions, while the oracle group miss shows why the two goals remain distinct. For zeros, calibrate a boundary-compatible score on the original response or declare a transformed/inverse-image target explicitly. No single raw or log-ratio geometry is an application-independent efficiency winner.

The synthetic design offers known ground truth, controlled heterogeneity and repeatable independent calibration experiments. It lacks a real application, an application-specific loss for region size, and a diversity of fitted predictors. The generator's log-ratio mean belongs to the fitted basis. Observed groups are supplied without estimation error, and the shift preserves their conditional laws exactly. Unknown groups, misspecified density ratios, changes within groups and deployment-time dependence can invalidate the corresponding interpretation. Neither a learned scale nor good empirical subgroup coverage establishes distribution-free pointwise conditional coverage.

Euclidean simplex area is only one reference measure. It emphasizes different parts of the simplex from a log-ratio measure and is evaluated only at $K=3$. Results at $K=10$ concern coverage and declared within-geometry diagnostics. We do not infer a high-dimensional efficiency ranking from low-dimensional area. Pointwise intervals across many cells do not provide multiplicity-controlled discoveries or prove exact validity from finite simulation.

Ordinary clr is undefined at zeros. A positive replacement applied consistently to all splits produces a different, exchangeable response $T_\epsilon(Y)$; its coverage does not guarantee direct membership of the original response in the transformed clr region. An explicit inverse-image region $\{y:T_\epsilon(y)\in C_\epsilon(x)\}$ does inherit the same coverage for original compositions, including zeros, by exact event identity. Its original-simplex area was not evaluated. Censoring in the sensitivity study is a measurement mechanism, not a model for every kind of structural zero. Choosing a scientifically justified zero mechanism remains application dependent.

We compare split-conformal score and calibration choices with one shared predictor. Full conformal, jackknife+, learned density regions and other multivariate methods may trade computation, adaptation and efficiency differently. Their absence restricts the comparison; it does not imply inferiority. No novelty claim rests on exhaustive discovery of the literature, and no conclusion is claimed for latent class probabilities inferred from one observed label.

\paragraph{Potential impacts.}
Better diagnostics can prevent users from treating an average coverage guarantee as assurance for an undercovered population. Conversely, presenting a compact region or a nominal guarantee without its target, group and shift conditions can encourage unwarranted confidence in consequential decisions. This synthetic study neither measures downstream harm nor certifies equitable decisions. Deployment would require relevant data, subgroup definitions, scientific treatment of zeros and monitoring for changes beyond the tested mixture intervention.
